\documentclass[11pt]{article}
\usepackage{amsmath,amssymb,amsthm,mathtools}
\usepackage[margin=1in]{geometry}
\usepackage{hyperref,xurl}
\setlength{\emergencystretch}{3em}
\newcommand{\C}{\mathbb C}
\newcommand{\R}{\mathbb R}
\newcommand{\im}{\operatorname{im}}
\newcommand{\diag}{\operatorname{diag}}
\title{Exact original-period metric conversion and the full factorial source return}
\author{}
\date{20 September 2026}
\begin{document}
\maketitle

\section{The original spaces and the two output norms}
Retain the original simple quartet, the fixed period and branch, the
complete arithmetic unit, and the actual source density
\[
 w_h(y)=\frac{|(g/h)(1/2+iy)|^2}{2\pi},\quad g=2\xi,
 \quad m_k=w_h^{*k},\quad c=k/2,\quad q=(k+1)^2.                  \tag{PM1}
\]
The physical polynomial norm is
$\int_\R|P(c+iy)|^2m_k(y)\,dy$.  The original full root polynomial
has distinct roots
\[
 \lambda_{ab}=c+(2a-k)\delta+i(2b-k)\gamma,
 \quad0\le a,b\le k,
 \qquad E_k=\C[S]/\chi_k.                                      \tag{PM2}
\]
For every original cutoff $N\ge q-1$, denote by $G_N$ the attained
metric on $E_k$, including all its polynomial relations and the
minimum over every degree-$N$ representative.  Let
$\Lambda:E_k\to B_k$ be the original observation onto its actual
image, in any one fixed coordinate basis of $B_k$, and retain
\[
 Q_N=(\Lambda G_N^{-1}\Lambda^*)^{-1},\qquad
 L_N=G_N^{-1}\Lambda^*Q_N.                                     \tag{PM3}
\]
All inverses exist since $\Lambda$ is onto its image and $G_N$
is positive definite.  Direct multiplication proves
$\Lambda L_N=I$, $L_N^*G_NL_N=Q_N$ and
$x=L_N\Lambda x+n$ with $n\in\ker\Lambda$ orthogonal in $G_N$
to the first summand.  Thus $Q_N$ is exactly the full observation
minimum.

The inherited period norm in the incoming
\path{01_CONDUCTOR_SUPPORT_AND_128_OBSERVATIONS.md}
uses the specified injection $\jmath:B_k\hookrightarrow\mathcal H_{\rm per}$.
Write $F_0$ for its original word realization, $G_0=F_0^*F_0$ for
the complete original word Gram, and $J$ for the inclusion of
the chosen $B_k$ coordinates into the word-coordinate space.
Then
\[
 \jmath=F_0J,\qquad P=J^*G_0J=\jmath^*\jmath>0,
 \qquad\jmath\Lambda=F_0\mathsf B\mathsf V_k.                   \tag{PM4}
\]
Here $\mathsf V_k$ is the unchanged original root-value map and
$\mathsf B$ the actual coefficient observation.  PM4 includes
every original word multiplicity, period entry and unit.  It
constructs $P$ rather than declaring it to be $Q_N$.

The exact comparison matrix and constants are
\[
 \Gamma_N=Q_N^{-1/2}P Q_N^{-1/2},\qquad
 \rho_N=\sqrt{\lambda_{\max}(\Gamma_N)},\qquad
 \underline\rho_N=\sqrt{\lambda_{\min}(\Gamma_N)}>0.             \tag{PM5}
\]
Substitution $v=Q_N^{1/2}b$ proves both sharp inequalities
\[
 \underline\rho_N\|b\|_{Q_N}\le\|b\|_P
                         \le\rho_N\|b\|_{Q_N}.                 \tag{PM6}
\]
The sharp upper constant has the original-map identity
\[
 \boxed{\rho_N=\|\jmath\Lambda\|_{(E_k,G_N)\to\mathcal H_{\rm per}}.}
                                                                    \tag{PM7}
\]
Indeed for every source vector the attained decomposition gives
$\|\Lambda x\|_{Q_N}\le\|x\|_{G_N}$, so PM6 proves one
inequality.  Conversely take $x=L_Nb$ for a maximizing vector of
PM5; its norm equals $\|b\|_{Q_N}$, giving equality.  The same
constant applies to a direct sum of any number of output blocks.
These are exact finite metric identities, not estimates inferred
from the number of observation steps.

\section{The unchanged root-value inverse and its metric receiver}
Let $M=M_S$ and let the actual finite recovery map in the incoming
support proof be
\[
 F_L:B_k^{\oplus(L+1)}\longrightarrow\C^q,
 \quad F_L(\Lambda x,\Lambda Mx,\ldots,\Lambda M^Lx)
                                  =\mathsf V_kx.               \tag{PM8}
\]
Its coefficients are the original translated conductor Lagrange
filters for the covered coordinates, followed by subtraction of
those values and the original detected-corner Lagrange filters.
No determinant estimate is substituted for these filters.
In particular $L\le127$ in that note, whereas the separate
constant-depth corner-block proof supplies a different valid map
with $L\le80$.  Either actual matrix can be used in what follows.

The exact norm in the inherited period metric is
\[
 C_{P,L}=\|F_L\diag(P^{-1/2},\ldots,P^{-1/2})\|_{2\to2}.       \tag{PM9}
\]
The O5 expression in the incoming proof is a bound for this
constant, in this precise metric.  With its displayed covered
coordinate bound $R_{\rm cov}$, remaining-corner bound $R_C$
and forward-stack bound $S_k$, the algebraic estimate is
\[
 C_{P,L}\le R_{\rm cov}+R_C(1+S_kR_{\rm cov}).                  \tag{PM10}
\]
For completeness, this last propagation is a direct operator
identity: if $F_{\rm cov}$ returns the covered coordinates,
$S_{\rm cov}$ inserts their forward observation stack, and
$F_C$ returns the remaining corner coordinates after subtraction,
the full map is
$F_{\rm cov}+F_C(I-S_{\rm cov}F_{\rm cov})$ in disjoint
root-coordinate blocks.  The triangle inequality and the three
stated norms give PM10.  No change to the underlying output
metric is involved in that estimate.  The factorial arithmetic
return below is independent of the word-Gram estimates used to
bound those three period-norm constants.

For a stack measured instead in $Q_N$, PM6 gives the exact
receiving inequality
\[
 \boxed{\|F_L\|_{Q_N^{\oplus(L+1)}\to\ell^2}
                      \le C_{P,L}\rho_N.}                     \tag{PM11}
\]
The factor $\rho_N$ is essential unless its actual matrix is
evaluated or bounded separately.  Equality of the kernels of
$\Lambda$ and $\jmath\Lambda$ gives no permission to erase it.
Conversely, a native-output bound transfers to the inherited
period input norm with factor $\underline\rho_N^{-1}$.

\section{Full original Lagrange lift and factorial convolution bound}
Let
\[
 d_*=2\min(\delta,\gamma),\qquad
 R_h=k\sqrt{\delta^2+\gamma^2},\qquad
 \eta=\pi/4,\qquad n=q-1.                                     \tag{PM12}
\]
Every original grid difference has modulus at least $d_*$:
for two different index pairs it is
$2\delta(a-a')+2i\gamma(b-b')$, and at least one integer
difference is nonzero.  The exact original cardinal polynomial
for a root $\lambda$ is
\[
 e_\lambda(S)=\prod_{\mu\ne\lambda}
                         \frac{S-\mu}{\lambda-\mu}.
\]
At the physical argument $S=c+iy$, every numerator factor is
bounded by $|y|+R_h$, since
$|\mu-c|\le k\sqrt{\delta^2+\gamma^2}$.  Hence, without changing
the full polynomial or coordinate,
\[
 |e_\lambda(c+iy)|^2
              \le d_*^{-2n}(|y|+R_h)^{2n}.                     \tag{PM13}
\]
For values $v=(v_\lambda)$, the literal Lagrange lift is
$P_v=\sum_\lambda v_\lambda e_\lambda$.  Cauchy--Schwarz in
this finite sum yields
\[
 |P_v(c+iy)|^2\le q\,d_*^{-2n}(|y|+R_h)^{2n}\|v\|_2^2.       \tag{PM14}
\]

The fixed full quartet polynomial $h$ has real coefficients,
because its roots and multiplicities are closed under conjugation.
The original entire function $g=2\xi$ is of real type.  Thus the
physical density in PM1 is even, and so is $m_k$.  Define its
actual bilateral moment
\[
 M_h(t)=\int_\R e^{ty}w_h(y)\,dy.
\]
The retained original source estimates prove
$M_h(\pi/2)<\infty$ (CAI16's $M_\alpha$).  Evenness implies
$M_h(t)=\int\cosh(ty)w_h(y)\,dy$, so in particular
$0<M_h(\eta)\le M_h(\pi/2)<\infty$.
Tonelli applied to the nonnegative convolution integrand gives
the exact unnormalized convolution identity
\[
 \int_\R e^{ty}m_k(y)\,dy=M_h(t)^k\quad(t=\eta,-\eta).
                                                                    \tag{PM15}
\]
Consequently
\[
 \int_\R e^{\eta|y|}m_k(y)\,dy
 \le\int_\R(e^{\eta y}+e^{-\eta y})m_k(y)\,dy
             =2M_h(\eta)^k.                                   \tag{PM16}
\]
This bound retains both half-lines and the complete original
convolution mass.  No probability normalization has occurred.

For $t\ge0$, the single $(2n)$th term of the exponential series
gives $t^{2n}\le(2n)!\eta^{-2n}e^{\eta t}$.  Insert
$t=|y|+R_h$ into PM14 and use PM16.  The original polynomial
norm therefore satisfies the fully evaluated bound
\[
 \boxed{\|P_v\|_{m_k}^2\le B_{h,k}\|v\|_2^2,\qquad
 B_{h,k}=2q(2q-2)!(\eta d_*)^{-2q+2}
                    e^{\eta R_h}M_h(\eta)^k.}                 \tag{PM17}
\]
This is exactly the incoming O6 constant, with its factor two,
factorial, physical displacement and mass all accounted for.
Since $P_v$ has degree $q-1\le N$, it is an admissible original
representative of its class at every cutoff under consideration.
Taking the full minimum only lowers its squared norm.  Hence
\[
 \|\mathsf V_k^{-1}v\|_{G_N}^2\le B_{h,k}\|v\|_2^2
                           \quad(N\ge q-1).                   \tag{PM18}
\]

For fixed original $h,\delta,\gamma$, the logarithm of the
right-hand constant is $O_h(q\log(q+2))$.  A direct sufficient
bound uses $(2q-2)!\le(2q-2)^{2q-2}$; the remaining factors in
PM17 have logarithms bounded in modulus by a constant times
$q+k+\log(q+2)$.  There is no altered asymptotic constant hidden
inside the proved exact formula PM17.

\section{The two recovery estimates and exact receiving matrix}
Combine PM9 and PM18, using the actual reconstruction
$\mathcal R_{L,N}=\mathsf V_k^{-1}F_L$.  It satisfies
\[
 \begin{split}
 \|\mathcal R_{L,N}\|_{P^{\oplus(L+1)}\to G_N}
      &\le\sqrt{B_{h,k}}\,C_{P,L},\\
 \|\mathcal R_{L,N}\|_{Q_N^{\oplus(L+1)}\to G_N}
      &\le\sqrt{B_{h,k}}\,C_{P,L}\rho_N.
 \end{split}                                                     \tag{PM19}
\]
The first is the proper original-metric interpretation of O6,
and is uniform over $N\ge q-1$ in its stated input norm.  Its
incoming $C_{P,L}=\exp(O_{h,u}(k))$ estimate gives the stated
$\exp(O_{h,u}(q\log(q+2)))$ return in that norm.  The second
estimate is the actual transfer to the attained output metric
used by DO and DQ.  Its factor $\rho_N$ has been specified and
proved in PM5--7; neither that factor nor its dependence on the
cutoff is estimated by the algebraic observation depth.

The full operator norm can also be evaluated directly, rather
than through these successive upper bounds.  Set
\[
 G_{{\rm val},N}=\mathsf V_k^{-*}G_N\mathsf V_k^{-1},\qquad
 Q_N^{[L]}=\diag(Q_N,\ldots,Q_N).
\]
Then the exact native recovery constant is
\[
 \boxed{\|\mathcal R_{L,N}\|^2
 =\lambda_{\max}\left((Q_N^{[L]})^{-1/2}
          F_L^*G_{{\rm val},N}F_L(Q_N^{[L]})^{-1/2}\right).}     \tag{PM20}
\]
The identity follows by substituting the explicit coordinate
isometries in the Rayleigh quotient, without changing any
original Gram.  It supplies a direct finite receiver for the
quantitative leakage bounds if the actual coefficient matrix
is preferred to PM19.  It does not imply a small determinant
allowance or a signed current estimate.

\section{Proof inputs and what was independently checked}
The original support inverse and its inherited output realization
are the actual O1--6 in the user-supplied
\path{01_CONDUCTOR_SUPPORT_AND_128_OBSERVATIONS.md};
the older corner-jet construction O1--12 is retained in its
accompanying input source.  The density and convolution mass
are the original LR2 and CAI16 definitions, including $1/(2\pi)$.
PM5--7 and PM11 prove the previously unstated metric conversion;
PM12--18 give a complete direct derivation of the original
factorial convolution bound.  The finite word-Gram and jet
estimates entering the incoming expression PM10 are kept as
their supplied period-norm bounds, not reclassified as native
quotient estimates.  All independent finite-dimensional and
integral arguments invoked here are included in the text.
\end{document}
